\section{Постановка задачи}
\label{sec:task}

Целью работы является исследование безопасности программного кода при обработке конфиденциальных и неконфиденциальных данных в оперативной памяти, а также построение методики сравнительной оценки реализаций.

В части~1 необходимо разработать приложение с интерактивным интерфейсом, позволяющее создавать, обновлять и удалять данные непосредственно в ОЗУ. После ввода конфиденциальных значений требуется выполнить их шифрование или хеширование. Далее следует проанализировать защищённость кода, использование оперативной памяти и процессорного времени, снять дампы памяти после создания, обновления и удаления записей и исследовать следы хранения данных.

В части~2 необходимо разработать методику оценки безопасности кода, применить её к собственной реализации и к предоставленному алгоритму, сравнить варианты и сформулировать выводы.
